\chapter{Incidencia dodecafásica y corredor Witt--Leech--Monstruo}
\label{chap:bandera-excepcional-acumulativa}
\index{doble proyección!bandera común}
\index{Witt!diseño}
\index{Monstruo}

El \cref{chap:paley-codigo-witt} ha construido explícitamente la matriz
\(A_W\), el código \(C_W\) y el sistema \(W_{12}\). Sobre esas coordenadas,
la doble proyección se formaliza mediante dos aplicaciones de evaluación que,
en el calibre declarado, producen por operaciones distintas la misma bandera
polarizada de Witt. Esta bandera calibrada determina además un origen marcado. La
construcción alcanza el diseño \(W_{12}\), su grupo de automorfismos
\(M_{12}\) y, mediante una segunda rama que no pasa por \(W_{24}\), un
3-vecino par unimodular sin raíces. La clasificación clásica identifica este
último con el retículo de Leech; la extensión al álgebra de operadores de
vértice lunar y al grupo Monstruo se obtiene mediante los teoremas clásicos
correspondientes.

\section{Estado dodecafásico y soportes de incidencia}

El vector dodecafásico es
\[
K=(234,543,140,729,659,824,621,058,914,794,146,601).
\]
Las cuatro palabras del marco son
\[
w_\pi=010211,
\quad
w_e=201101,
\quad
w_\varphi=121200,
\quad
w_{\rm APP}=022020.
\]
En el calibre simétrico de Paley--Witt se usa el codificador
\[
c(w)=(w,wA_W)\in\mathbb F_3^{12}.
\]
Sus soportes son
\[
\begin{aligned}
P_\pi&=\{2,4,5,6,7,8,9,10,11\},\\
H_e&=\{1,3,4,6,9,10\},\\
H_\varphi&=\{1,2,3,4,7,9\},\\
H_{\rm APP}&=\{2,3,5,10,11,12\}.
\end{aligned}
\]
\(P_\pi\) tiene nueve puntos. Los otros tres soportes son hexadas. Esta
distinción debe conservarse: las doce hexadas contenidas en \(P_\pi\) poseen
como triples complementarios las rectas de \(AG(2,3)\).

\section{Coincidencia de los soportes de umbral}

La evaluación arquimediana extrae del vector \(K\) el soporte de umbral
\[
\beta_{\rm arq}(K)=\{i:K_i\ge729\}
=\{4,6,9,10\}.
\]
La evaluación excepcional calcula la intersección
\[
\beta_{\rm exc}(w_e,w_\pi)
=\operatorname{supp}c(w_e)\cap\operatorname{supp}c(w_\pi).
\]
La evaluación produce la primera igualdad exacta condicionada a estos datos:
\begin{equation}
\boxed{
\beta_{\rm arq}(K)
=\beta_{\rm exc}(w_e,w_\pi)
=B_K=\{4,6,9,10\}.}
\label{eq:bandera-cara-alta}
\end{equation}
Las dos aplicaciones realizan operaciones distintas sobre entradas
declaradas: la primera usa umbrales del vector \(K\), mientras que la segunda
usa intersecciones de soportes del código. El umbral \(729\), las palabras y
el calibre de Paley forman parte de la calibración. En particular, esta
igualdad no es una validación independiente de \(K\): todos los umbrales
enteros entre \(660\) y \(729\), inclusive, seleccionan el mismo soporte
\(B_K\).

\section{Coincidencia de los soportes negativos}

Antes del acarreo, la cocadena arquimediana es
\[
s^{(0)}=P+E-\Phi-K,
\]
y sus doce componentes son
\[
(7,297,353,-430,-715,-1199,-3,286,-894,-528,380,663).
\]
Su soporte negativo vale
\[
\eta_{\rm arq}(s^{(0)})
=\{4,5,6,7,9,10\}.
\]
Por la ruta excepcional, fijamos como selector de calibración la hexada \(H\)
que satisfaga
\[
B_K\subset H\subset P_\pi
\]
y
\[
(|H\cap H_e|,|H\cap H_\varphi|,|H\cap H_{\rm APP}|)=(4,3,2).
\]
La enumeración del diseño produce una única solución para ese selector:
\[
\eta_{\rm exc}
=\{4,5,6,7,9,10\}
=H^-_\alpha.
\]
Por tanto,
\begin{equation}
\boxed{
\eta_{\rm arq}(s^{(0)})
=\eta_{\rm exc}(B_K,P_\pi,H_e,H_\varphi,H_{\rm APP})
=H^-_\alpha.}
\label{eq:bandera-hexada-negativa}
\end{equation}

Las ecuaciones~\eqref{eq:bandera-cara-alta} y
\eqref{eq:bandera-hexada-negativa} se condensan en la bandera
\[
\mathcal I_*^{\rm core}=(B_K\subset H^-_\alpha).
\]
Ésta es la bandera de compatibilidad calibrada de la doble proyección. El marco marcado
\[
\widetilde{\mathcal I}_*
=(\mathcal I_*^{\rm core};P_\pi,H_e,H_\varphi,H_{\rm APP})
\]
conserva los datos necesarios para construir la rama excepcional. Bajo una
permutación simultánea \(g\in M_{12}\), toda la construcción se transporta
equivariantemente; lo invariante es su clase de isomorfía, no los nombres
numéricos aislados.

\section{Reconstrucción del punto base a partir del marco de incidencia}

El origen no necesita quedar como selector exterior si se conserva el marco
ordenado completo. Definimos
\[
o(\mathcal F)
=(H_e\cap H_\varphi)\setminus(H^-\cup H_{\rm APP}),
\qquad
\mathcal F=(H^-,H_e,H_\varphi,H_{\rm APP}).
\]

\begin{theorem}[Origen determinado por incidencias]
Para el marco canónico,
\[
o(\mathcal F_*)=\{1\}.
\]
Además, la regla es equivariante: \(o(g\mathcal F)=g\,o(\mathcal F)\).
\end{theorem}

\begin{proof}
Se tiene
\[
H_e\cap H_\varphi=\{1,3,4,9\},
\]
y
\[
(H^-_\alpha\cup H_{\rm APP})\cap\{1,3,4,9\}
=\{3,4,9\}.
\]
La diferencia es \(\{1\}\). La equivarianza sigue de que intersección,
unión y diferencia conmutan con la acción de cualquier permutación.
\end{proof}

El paso desde este origen de incidencia hasta un vector de un retículo
\(A_2^{12}\) requiere fijar una realización métrica de las doce coordenadas.
El \cref{sec:vecino-leech-ejecutable} fija esa realización, muestra
que el origen selecciona uno de los doce minimizadores radiales y certifica
las condiciones del 3-vecino. La incidencia determina el origen; la matriz de
Gram y la cámara positiva determinan después el vector.

\section{Incidencia del diseño ternario}

La cadena de Paley
\[
C_P\longmapsto A_W\longmapsto C_W\longmapsto W_{12}
\]
se encuentra fijada por las matrices y demostraciones del
\cref{chap:paley-codigo-witt}. En particular,
\(A_W^2=-I_6\), \(C_W=[12,6,6]_3\), y los 132 soportes proyectivos de
peso seis forman \(W_{12}=S(5,6,12)\). El paso de una palabra a su soporte
identifica exactamente el par \(\{c,-c\}\); las incidencias empleadas en lo
que sigue pertenecen, por consiguiente, al diseño proyectivo y no a palabras
signadas individuales.

\section{Hexadas, octadas y la identidad de cardinalidad 28}

El diseño binario \(W_{24}=S(5,8,24)\) posee octadas, cada una con
\(\binom82=28\) pares. El \cref{thm:dodecada-w12} construye
\(W_{12}\) como diseño residual de una dodecada \(D\) del código binario
extendido de Golay; el \cref{thm:elevacion-unica-hexada} demuestra que, una
vez fijado el alineamiento \(\ell\), cada hexada posee una elevación única
\(O_H\) a una octada. El paso \(6\to8\) es, por tanto, un morfismo de
incidencia relativo al dato binario \((D,\ell)\), y no una inferencia basada
en cardinalidades.

Por otro lado, la microfibra de \(\pi\) contiene 1008 rutas y una nónada tiene
36 pares, de modo que
\[
1008=36\cdot28.
\]
La identidad \(1008=36\cdot28\) es un control numérico del catálogo. El
morfismo de incidencia se define independientemente mediante
\(H\mapsto\iota_D(\ell(H))\); la identidad de cardinalidades no interviene en
su definición.

\section{Involuciones de estrella y generación del grupo de Mathieu}

Para cada cara \(B\in\binom{[12]}4\), existen cuatro hexadas de \(W_{12}\)
que contienen \(B\). Los cuatro residuos \(H\setminus B\) son pares disjuntos
que particionan \([12]\setminus B\). La involución \(\sigma_B\) fija \(B\)
e intercambia los extremos de esos cuatro pares. Hay
\[
\binom{12}{4}=495
\]
involuciones de estrella.

La enumeración exhaustiva, para la matriz y las palabras declaradas, verifica
\[
\sigma_B(W_{12})=W_{12}
\]
para todas ellas y
\[
\left|\langle\sigma_B:B\in\tbinom{[12]}4\rangle\right|=95040.
\]
Como \(\Aut(W_{12})\cong M_{12}\), las 495 estrellas generan \(M_{12}\).
La cara HMT es
\[
B_K=\{4,6,9,10\},
\qquad
\sigma_{B_K}=(1\ 3)(2\ 11)(5\ 7)(8\ 12).
\]
Una sola estrella genera sólo \(C_2\). La acción combinatoria global se
expresa mediante la representación del grupo libre de generadores de estrella
\[
\rho_\star:F_{495}\longrightarrow\Aut(W_{12}),
\qquad
\rho_\star(\gamma_B)=\sigma_B,
\qquad
\operatorname{im}\rho_\star=M_{12},
\]
donde \(F_{495}\) es el grupo libre sobre las 495 estrellas formales. El grupo
global es, por tanto, la imagen de esa familia de involuciones y no la de una
sola estrella. Identificar estos generadores con lazos del espacio de estados
requeriría construir, además, el transporte correspondiente; esa
identificación no interviene en el cálculo del grupo.

\section{Dos prolongaciones excepcionales del código ternario}

Desde el código ternario pueden continuarse dos ramas de tipo distinto:
\[
\begin{aligned}
C_W&\longrightarrow W_{12}
\xrightarrow[\ (D,\ell)\ ]{\ \iota_D\ }W_{24},\\
C_W&\xrightarrow[\phi]{\text{pegado explícito}}N(A_2^{12})
\xrightarrow[v_\alpha]{\text{3-vecino exacto}}\Lambda_{24}
\longrightarrow V^\natural .
\end{aligned}
\]
La primera rama continúa mediante el diseño residual de la dodecada y su
elevación única; la segunda continúa mediante un pegado ternario y un vecino
de índice tres. Las dos construcciones tienen codominios distintos. Para la
segunda, sea \(A_2=\mathbb Z\alpha_1\oplus\mathbb Z\alpha_2\) con matriz
de Gram
\[
G_{A_2}=\begin{pmatrix}2&-1\\-1&2\end{pmatrix},
\qquad \rho=\alpha_1+\alpha_2,
\]
y tomemos como representantes de \(A_2^*/A_2\cong\mathbb F_3\)
\[
\varpi_0=0,\qquad
\varpi_1=\frac{2\alpha_1+\alpha_2}{3},\qquad
\varpi_2=\frac{\alpha_1+2\alpha_2}{3}.
\]
La realización de pegado es la aplicación coordenada a coordenada
\[
\phi(c_1,\ldots,c_{12})=(\varpi_{c_1},\ldots,\varpi_{c_{12}}),
\]
que forma
\[
N(C_W)=\bigcup_{c\in C_W}(A_2^{12}+\phi(c)).
\]
La linealidad de \(C_W\) hace que la unión de clases laterales sea un retículo; su
autodualidad da la integralidad del pegado, y la divisibilidad por tres de
todos los pesos da su paridad. Además, \(C_W\) tiene dimensión seis y
\[
\det N(C_W)=\frac{3^{12}}{|C_W|^2}
=\frac{3^{12}}{(3^6)^2}=1.
\]
La distancia mínima seis implica que una clase de pegado no nula tiene norma
al menos \(6(2/3)=4\). Por tanto las raíces de \(N(C_W)\) son exactamente las
de \(A_2^{12}\). Ésta es la construcción explícita del Niemeier de tipo
\(A_2^{12}\) usada en lo que sigue.

\section{Minimización radial y construcción explícita del 3-vecino}
\label{sec:vecino-leech-ejecutable}

En la cámara positiva consideramos vectores radiales
\[
v(a)=\sum_{i=1}^{12}a_i\rho_i,
\qquad a_i=1+3b_i,\quad b_i\ge0.
\]
La condición \(v(a)^2\equiv0\pmod{18}\) equivale a
\(\sum_i b_i\equiv1\pmod3\). Su mínimo es
\[
v(a)^2=2\sum_i(1+3b_i)^2=54,
\]
y se alcanza exactamente doce veces: un coeficiente vale cuatro y los once
restantes valen uno. El origen recuperado \(o_*=1\) selecciona
\[
v_\alpha=4\rho_1+\rho_2+\cdots+\rho_{12}.
\]
La minimalidad pertenece al cono radial positivo declarado. En todo el
cociente residual existe, por ejemplo, el representante
\((\alpha_1-2\alpha_2,\rho,\ldots,\rho)\), congruente con
\((\rho,\ldots,\rho)\) módulo tres y de norma \(14+11\cdot2=36\). El
cálculo directo lo conserva como testigo que impide ampliar el dominio.

Definimos
\[
N_{\alpha,0}=\{x\in N(C_W):
\langle x,v_\alpha\rangle\equiv0\pmod3\},
\qquad
N_\alpha=N_{\alpha,0}+\mathbb Z\frac{v_\alpha}{3}.
\]

\begin{theorem}[3-vecino de Leech]
La construcción anterior satisface
\begin{align*}
\mathrm{(I1)}\quad&v_\alpha\in A_2^{12}\subset N(C_W),
&\langle N(C_W),v_\alpha\rangle&\not\subset3\mathbb Z,\\
\mathrm{(I2)}\quad&v_\alpha^2=54\equiv0\pmod{18},\\
\mathrm{(I3)}\quad&v_\alpha/3\in
N_{\alpha,0}^*\setminus N(C_W),
&3(v_\alpha/3)&\in N_{\alpha,0}.
\end{align*}
El retículo \(N_\alpha\) es par, unimodular, de rango veinticuatro y sin
raíces. En consecuencia,
\[
N_\alpha\cong\Lambda_{24}.
\]
\end{theorem}

\begin{proof}
Para cualquier raíz de un factor \(A_2\), el emparejamiento con
\(a_i\rho_i\) es congruente con \(1\) o \(2\) módulo tres. El carácter de
(I1) es, por ello, sobreyectivo, y ninguna de las raíces antiguas pertenece a
\(N_{\alpha,0}\). La igualdad de (I2) es el cálculo radial anterior. Si
\(x\in N_{\alpha,0}\), entonces
\(\langle x,v_\alpha/3\rangle\in\mathbb Z\), lo que demuestra la primera
inclusión de (I3). Si \(v_\alpha/3\) perteneciera a \(N(C_W)\), la
integralidad del Niemeier obligaría a que el carácter de (I1) fuese nulo. La
clase tiene orden tres y \(v_\alpha\in N_{\alpha,0}\) porque
\(v_\alpha^2\equiv0\pmod3\).

El núcleo tiene índice tres, luego
\[
\det N_{\alpha,0}=3^2\det N(C_W)=9,
\qquad
\det N_\alpha=9/3^2=1.
\]
Para \(x\in N_{\alpha,0}\) y \(m\in\mathbb Z\),
\[
\left(x+m\frac{v_\alpha}{3}\right)^2
=x^2+2m\frac{\langle x,v_\alpha\rangle}{3}+6m^2
\in2\mathbb Z,
\]
así que el vecino es par. Queda excluir raíces nuevas. Cada coordenada de
las dos clases laterales no triviales pertenece a una de las seis clases
\[
A_2+\varpi_j\pm\rho/3,\qquad j\in\{0,1,2\},
\]
y el mínimo de cada clase es \(2/9\). En efecto, al escribir una coordenada
como \((u/3,v/3)\), la norma es
\(2(u^2-uv+v^2)/9\); los seis residuos son no nulos y cada uno admite un
representante con \(u^2-uv+v^2=1\). Toda
norma nueva es al menos \(12(2/9)=8/3>2\). El vecino no tiene raíces. La
clasificación de los retículos pares unimodulares de rango veinticuatro lo
identifica con Leech.
\end{proof}

La enumeración de los seis mínimos locales y los doce minimizadores radiales
reconstruye las \(729\) palabras de \(C_W\), la matriz de Gram del pegado y
su paridad; la aritmética exacta comprueba I1--I3, los determinantes y la cota
\(8/3\).

\section{Corredor de Leech al álgebra Moonshine}
\label{sec:corredor-hmt-moonshine}

Definamos
\[
 \Lambda_{\rm HMT}:=N_\alpha\cong\Lambda_{24},
 \qquad
 V^\natural_{\rm HMT}
 :=V_{\Lambda_{\rm HMT}}^+
 \oplus V_{\Lambda_{\rm HMT}}^{T,+}.
\]
No se introduce aquí una segunda retícula: la construcción
Frenkel--Lepowsky--Meurman se aplica al vecino producido en el apartado
anterior.

\begin{theorem}[Corredor HMT--Moonshine]
\label{thm:corredor-hmt-moonshine}
Fijados la dodecada, el alineamiento, el origen y la cámara excepcional de la
construcción precedente, se cumplen
\[
 V^\natural_{\rm HMT}\cong V^\natural,
 \qquad
 \Aut(V^\natural_{\rm HMT})\cong\mathbb M,
\]
y
\[
 \operatorname{Tr}_{V^\natural_{\rm HMT}}q^{L_0-1}
 =J_{\rm Mon}(\tau)
 =j(\tau)-744
 =q^{-1}+196884q+21493760q^2+\cdots .
\]
En consecuencia, la cadena
\[
 \mathrm{HMT}\longrightarrow N_\alpha\cong\Lambda_{24}
 \longrightarrow V^\natural_{\rm HMT}
 \longrightarrow\mathbb M
 \longrightarrow\text{\rm Moonshine}
\]
es un resultado demostrado por composición.
\end{theorem}

\begin{proof}
El teorema del 3-vecino demuestra, a partir de los datos declarados, que
\(N_\alpha\) es par, unimodular, de rango veinticuatro y sin raíces. La
clasificación de retículos lo identifica con \(\Lambda_{24}\). Aplicada a
esta salida, la construcción de Frenkel--Lepowsky--Meurman produce el
orbifold
\(V_{\Lambda_{\rm HMT}}^+\oplus V_{\Lambda_{\rm HMT}}^{T,+}\);
los teoremas clásicos identifican su grupo de automorfismos con
\(\mathbb M\) y su carácter graduado con el Hauptmodul normalizado
\(J_{\rm Mon}=j-744\)
\cite{FLM1988,ConwayNorton1979,Borcherds1992}. Cada flecha entrega, por
tanto, un objeto del tipo requerido por la siguiente.\qedhere
\end{proof}

\begin{remark}[La componente trivial no es un ajuste]
La gradación satisface
\[
 V^\natural_{\rm HMT}
 =\bigoplus_{n\geq0}(V^\natural_{\rm HMT})_n,
 \qquad
 (V^\natural_{\rm HMT})_0=\CC\mathbf1,
 \qquad
 (V^\natural_{\rm HMT})_1=0.
\]
El término \(q^{-1}\) registra así la línea de vacío. En grado dos,
\[
 \dim(V^\natural_{\rm HMT})_2
 =196884
 =1+196883
 =196560+324
 =54(5\cdot729+1).
\]
La igualdad \(196884=1+196883\) separa la representación trivial y la menor
representación irreducible no trivial del Monstruo. Las otras dos igualdades
son descomposiciones enteras exactas; no sustituyen esta estructura graduada
ni se emplean para seleccionar retrospectivamente el corredor.
\end{remark}

La etapa finita hasta \(N_\alpha\) parte de una dodecada, un alineamiento, un
origen y una cámara declarados. El cálculo prueba que \(N_\alpha\) es par,
unimodular y sin raíces. Sobre esa salida se aplican los teoremas clásicos de
clasificación y orbifold: el orbifold es \(V^\natural\), su grupo de
automorfismos es \(\mathbb M\) y su carácter graduado es \(J_{\rm Mon}\).
La composición conserva así la procedencia de cada flecha sin reducir el
corredor a una coincidencia de coeficientes.

\section{Espacio cociente de 243 síndromes}

Perforar una coordenada de \(C_W\) produce
\[
G_{11}=[11,6,5]_3.
\]
La bola ternaria de radio dos tiene
\[
1+22+220=243
\]
elementos. Como
\[
|G_{11}|\cdot243=3^6\cdot3^5=3^{11},
\]
y la distancia mínima es cinco, esas bolas son disjuntas y cubren todo
\(\mathbb F_3^{11}\). Cada vector admite una descomposición única
\[
x=c+e,
\qquad
c\in G_{11},
\quad
\operatorname{wt}(e)\le2.
\]
El cociente lineal es
\[
\mathbb F_3^{11}/G_{11}\cong\mathbb F_3^5,
\qquad
|\mathbb F_3^{11}/G_{11}|=243.
\]
Por tanto, se obtiene un espacio cociente de 243 síndromes en un ambiente de once
coordenadas y dimensión cociente cinco. El catálogo de 243 palabras \(w_6\)
y este cociente tienen el mismo cardinal, pero son objetos de tipo distinto;
el codificador directo lleva las palabras del código al síndrome cero. La doble proyección
demostrada utiliza la bandera común, no una biyección de cardinalidades sin
estructura.

\section{Compatibilidad finita de las evaluaciones aritmética y excepcional}

\begin{theorem}[Compatibilidad finita condicionada de las dos proyecciones]
Considérense el vector \(K\) y el estado dodecafásico firmado construidos en
el \cref{chap:alpha}; las cuatro palabras del catálogo definidas en el
\cref{chap:regiones}; y la matriz \(A_W\) construida en el
\cref{chap:paley-codigo-witt}. Fijada una calibración
\(c_{\rm exc}\in\mathscr C_{\rm exc}\), se cumplen las afirmaciones
siguientes:
\begin{enumerate}
  \item la evaluación arquimediana obtiene \(\alpha_{12}\) mediante el acarreo
  dodecafásico;
  \item los mapas de soporte arquimediano y excepcional producen la misma bandera
  \(B_K\subset H^-_\alpha\);
  \item el marco marcado recupera el origen \(o_*=1\);
  \item la matriz \(A_W\) construye \(W_{12}\), y las 495 involuciones de
  estrella generan un grupo de orden \(95040\), identificado con \(M_{12}\);
  \item la realización discriminante explícita construye \(N(A_2^{12})\),
  y el origen, junto con la cámara positiva, selecciona \(v_\alpha\); su
  3-vecino es par, unimodular y sin raíces.
\end{enumerate}
Las identidades de soportes y la regla del origen son equivariantes bajo una
permutación simultánea de las doce coordenadas. La calibración transforma de
manera simultánea los datos de incidencia y no altera su clase de isomorfía.
\end{theorem}

\begin{proof}
La primera afirmación es la reconstrucción dodecafásica del
\cref{chap:alpha}. Las ecuaciones
\eqref{eq:bandera-cara-alta} y \eqref{eq:bandera-hexada-negativa} demuestran
la segunda por evaluación directa de las dos aplicaciones. El cálculo conjuntista
de la sección anterior da \(o_*=1\). Finalmente, el enumerador, los 132
soportes de peso seis, las 495 involuciones y el orden del grupo, junto con el
teorema del 3-vecino, prueban la última afirmación.
\end{proof}

\begin{corollary}[Compatibilidad con el corredor HMT--Moonshine]
Los datos finitos de la proyección excepcional construyen el mismo
\(N_\alpha\) que interviene en el
\cref{thm:corredor-hmt-moonshine}; por tanto, la evaluación finita es
compatible con \(V^\natural_{\rm HMT}\), su acción de \(\mathbb M\) y su
carácter \(J_{\rm Mon}\).
\end{corollary}

\enlargethispage{3\baselineskip}
La bandera finita se obtiene mediante dos operaciones distintas aplicadas a
entradas calibradas, y el origen se deriva después de fijar esa bandera. Esto
prueba una identidad condicionada, no una procedencia independiente de las
entradas. La construcción reticular conserva la diferencia de tipo. La
composición desde los datos de \(\alpha\) hasta \(\Lambda_{24}\) factoriza por
el código, el pegado \(A_2^{12}\), el origen marcado, la cámara radial y el
3-vecino. La rama \(W_{12}\to W_{24}\) continúa separada y queda determinada
para toda elección de una dodecada \(D\) y un isomorfismo de incidencia
\(\ell:W_{12}^{\rm HMT}\to\mathcal H(D)\). El teorema es paramétrico en esos
datos y no fija una elección coordenada particular de \((D,\ell)\).
